\section{问题三建模与求解：通信约束下的运输与中继协同调度}

\subsection{三维空间射线追踪视距遮挡与双向链路预算模型}

\subsubsection{高精 DEM 空间射线追踪与视距（LOS）判据}
在山区复杂地形中，微波信号沿视线直线传播，易受沿途山体阻挡。
设发射端点 $a$ 与接收端点 $b$ 在三维大地空间坐标系下的三维坐标分别为 $(\text{lon}_a, \text{lat}_a, z_a)$ 与 $(\text{lon}_b, \text{lat}_b, z_b)$，空间欧氏直线距离为 $D_{ab}\text{ (km)}$。
沿连线方向以采样步长 $\Delta s = 20\text{ m}$ 离散取样，参数化表示为：
\begin{equation}
    \begin{cases}
        \text{lon}(\alpha) = (1 - \alpha) \text{lon}_a + \alpha \text{lon}_b \\
        \text{lat}(\alpha) = (1 - \alpha) \text{lat}_a + \alpha \text{lat}_b \\
        z_{LOS}(\alpha) = (1 - \alpha) z_a + \alpha z_b
    \end{cases}, \quad \alpha \in (0, 1)
    \label{eq:ray_tracing}
\end{equation}
定义地形视距遮挡指示变量 $b_{ab} \in \{0, 1\}$。若连线上存在任意插值点的视线海拔低于地面 DEM 高程，则判定为视距遮挡（NLOS）：
\begin{equation}
    b_{ab} = \begin{cases}
        1, & \exists \alpha \in (0, 1) \text{ s.t. } z_{LOS}(\alpha) \le z_{dem}\big(\text{lon}(\alpha), \text{lat}(\alpha)\big) \\
        0, & \text{otherwise (LOS)}
    \end{cases}
    \label{eq:nlos_indicator}
\end{equation}

\subsubsection{第一菲涅尔区保活几何机制与双向链路门限}
更严密地，无线电波在空间传播时存在有效能量包络。载波频率 $f = 2400\text{ MHz}$（波长 $\lambda = c / f = 0.125\text{ m}$）的第一菲涅尔区（First Fresnel Zone, $F_1$）在连线上任意点 $\alpha$ 的椭球面半短轴半径为：
\begin{equation}
    r_{F_1}(\alpha) = \sqrt{\frac{\lambda \cdot d_1 \cdot d_2}{d_1 + d_2}} = \sqrt{\frac{\lambda \cdot \alpha(1-\alpha) D_{ab}}{1}}
    \label{eq:fresnel_radius}
\end{equation}
若山脊不仅遮挡视线核心，而且刺入第一菲涅尔区，将诱发强烈的衍射损耗与电磁波相消干涉。
无线电全路径总损耗模型定义为自由空间扩散损耗（FSPL）与山体遮挡附加衰减之和：
\begin{equation}
    L_{path, ab} = 32.45 + 20 \log_{10}(f) + 20 \log_{10}(D_{ab}) + 10.0 \cdot b_{ab}
    \label{eq:path_loss}
\end{equation}
根据题面附录 3 标定的硬件灵敏度门限（接收灵敏度 $-98\text{ dBm}$，衰落裕量 $8\text{ dB}$，发射功率与天线增益匹配），建立三大关键通信链路的最大容忍损耗门限：
\begin{enumerate}
    \item \textbf{直连通信链路（$G01 \leftrightarrow UAV$）}：$L_{path} \le L_{max}^{G01 \leftrightarrow UAV} = 122.0\text{ dB}$；
    \item \textbf{中继回传链路（$Relay \leftrightarrow G01$）}：$L_{path} \le L_{max}^{Relay \leftrightarrow G01} = 126.0\text{ dB}$；
    \item \textbf{中继接入链路（$UAV \leftrightarrow Relay$）}：$L_{path} \le L_{max}^{UAV \leftrightarrow Relay} = 116.0\text{ dB}$。
\end{enumerate}

对问题二确定的 24 个运输架次开展时空射线追踪仿真。如图~\ref{fig:q3_fresnel}所示，以 O01 飞往 S003 的航线为例，沿途最高山脊（海拔 423.8m）不仅遮断直视视距，更刺入第一菲涅尔区深达 150 米以上，导致显著的附加衰落与信号阻断。
全空域仿真表明：全文 24 个运输航次共包含 325 处微时段，其中直连链路中断区间达 \textbf{136 处}（占比高达 \textbf{41.85\%}），必须依靠空中中继进行空间几何破障。

\begin{figure}[htbp]
  \centering
  \includegraphics[width=0.95\textwidth]{fig10_q3_los_dem_profile_blockage.pdf}
  \caption{30m DEM 山区复杂地形视距遮挡 (LOS) 射线追踪与第一菲涅尔区保活几何机制剖面}
  \label{fig:q3_fresnel}
\end{figure}

\subsection{DEM 山体战略制高点空中中继选址}

中继无人机 $R$ 必须悬停于特定空中制高点，充当天空中的“虚拟基站”。选址需同时满足两项刚性条件：
\begin{enumerate}
    \item \textbf{回传强保活}：与调度中心固定网关 G01 保持严格 LOS 通视，且回传损耗 $L \le 126.0\text{ dB}$；
    \item \textbf{前向广覆盖}：悬停高度满足题面离地高度上限（$\text{AGL} \le 300.0\text{ m}$），且与被遮挡运输机满足接入损耗 $L \le 116.0\text{ dB}$。
\end{enumerate}

通过对灾区 DEM 高程场 224 个山峰制高点进行三维网格遍历与视距分析，本文精确定位了三大战略悬停制高点（悬停离地高度均取上限 300.0m）：
\begin{itemize}
    \item \textbf{西区制高点 ($P_{west}$)}：坐标 $(109.20214^\circ\text{E}, 23.03643^\circ\text{N})$，地面海拔 431.6m，悬停绝对海拔 \textbf{731.6m}。与 G01 距离 4.21km，回传损耗仅 112.8 dB（远低于 126.0 dB 门限）。以高空俯视视野无死角覆盖 S003、S005、S006、S007、S008、S011、S015 等西部 7 个深谷受灾点；
    \item \textbf{东南制高点 ($P_{es}$)}：坐标 $(109.27857^\circ\text{E}, 23.03643^\circ\text{N})$，地面海拔 387.3m，悬停绝对海拔 \textbf{687.3m}。与 G01 回传损耗 115.3 dB，专向覆盖东南部深山孤立点 S014；
    \item \textbf{东北制高点 ($P_{en}$)}：坐标 $(109.27500^\circ\text{E}, 23.04786^\circ\text{N})$，地面海拔 376.9m，悬停绝对海拔 \textbf{676.9m}。与 G01 回传损耗 116.1 dB，有效覆盖东北部山脊背风区 S004 与 S002 的低空遮挡盲区。
\end{itemize}

\subsection{实体中继长续航轮转调度模型与方案对比}

中继无人机起飞总质量 $m_{takeoff} = 23.5\text{ kg}$，搭载专用的 3.2 kWh 锂能源组件。
往返飞行巡航功率 $P_{cruise} = 1.15\text{ kW}$；悬停建链阶段功率为 $1.05\text{ kW}$（30 秒），正式中继服务阶段综合功率为 $P_{hover} + P_{com} = 1.05 + 0.05 = 1.10\text{ kW}$。
调度中心仅配置 2 架实体中继机（R01, R02）与 6 组共享能源组件。
实体机返航后必须满足地面周转冷却与换电准备约束：$\Delta t_{turnaround} \ge 300.0\text{ s}$。

基于上述物理系统，本文构建了两种调度优化方案：
\begin{enumerate}
    \item \textbf{方案 A（3 架次精简分解方案，本文正式推荐）}：
    采用长周期静态驻留策略，仅安排 3 个长续航中继架次（RT01 $\sim$ RT03）。
    R01 执飞 RT01，长周期驻留西区制高点，贯穿保障全区主要航运；
    R02 执飞 RT02（东南区）与 RT03（东北区）。
    方案 A 的详细架次排程如表~\ref{tab:q3_schemeA}所示。
    
    \item \textbf{方案 B（4 架次跨区动态接力方案）}：
    在西区空档期安排中继机提前返航换电，拆分为 4 个较短架次实施跨区接力。
\end{enumerate}

\begin{table}[htbp]
  \centering
  \caption{推荐方案 A（3 架次精简分解）中继调度明细表}
  \label{tab:q3_schemeA}
  \fontsize{9pt}{12pt}\selectfont
  \begin{tabular}{ccccccccccc}
    \toprule
    \textbf{架次} & \textbf{实体机} & \textbf{组件} & \textbf{悬停制高点} & \textbf{海拔} & \textbf{起飞(s)} & \textbf{建链(s)} & \textbf{结束(s)} & \textbf{落地(s)} & \textbf{能耗(kWh)} & \textbf{返航SOC} \\
    \midrule
    \textbf{RT01} & R01 & MOD\_01 & 西区制高点 & 731.6m & 140.0 & 785.9 & 7260.0 & 7746.3 & 2.2227 & 30.54\% \\
    \textbf{RT02} & R02 & MOD\_02 & 东南制高点 & 687.3m & 0.0 & 735.7 & 2150.0 & 2722.3 & 0.7371 & 76.96\% \\
    \textbf{RT03} & R02 & MOD\_03 & 东北制高点 & 676.9m & 3025.0 & 3791.6 & 5250.0 & 5852.4 & 0.7711 & 75.90\% \\
    \bottomrule
  \end{tabular}
\end{table}

如图~\ref{fig:q3_timeline}所示，方案 A 的时序协同特征如下：
Panel (a) 清晰呈现了 R02 在完成 RT02（2722.3 秒落地）与起飞执行 RT03（3025.0 秒）之间的地面周转维护时间达到 $3025.0 - 2722.3 = \textbf{302.7 秒} \ge 300.0\text{ 秒}$，满足地面维护硬约束；
Panel (b) 展现了 24 个运输航次的 325 处通信微时段（直连 189 处占 58.15\%，中继 136 处占 41.85\%）在时间轴上实现 \textbf{100\% 连续覆盖}，全航程无通信中断盲区（相邻微时段重叠或间隔 $< 0.1\text{ s}$）。

\begin{figure}[htbp]
  \centering
  \includegraphics[width=0.95\textwidth]{fig11_q3_relay_schedule_timeline.pdf}
  \caption{通信中继无人机调度全景时序甘特图与 100\% 连续通信保障时空覆盖}
  \label{fig:q3_timeline}
\end{figure}

针对方案 A 与方案 B 的决策权衡，本文结合航次数量、起降频次、能源组件占用及全系统完工时效等指标进行综合对比（表~\ref{tab:q3_compare}与图~\ref{fig:q3_radar}）。
在全系统联合完工时间同为 7746.3 秒（由末架中继机 RT01 返航着陆时刻决定，运输机群已于 7730.5 秒前全数归巢）的前提下，两方案联合总能耗分别为 81.9963 kWh 与 81.8158 kWh，仅相差 0.1805 kWh（相对差异为 0.22\%）。方案 A 采用长效驻留策略，将中继架次由 4 架次精简为 3 架次，起降作业减少 1 次（降低 25\%）；同时仅占用 3 组能源组件（保留 3 组备用，应急储备率达 50\%），且避免了空中二次调机与链路交接过程中的通信中断隐患。方案 B 虽在悬停能耗上存在微小理论数值降低，但需要额外执行一次起降作业并多占用 1 组能源组件。在山地复杂气流与非确定性环境下，方案 A 以 0.22\% 的微小能耗差异换取了更低的起降频次、更高的组件储备以及连续稳定的通信链路，在现场操作与执行可靠性上具有明确的工程应用优势。

\begin{table}[htbp]
  \centering
  \caption{中继协同调度推荐方案 A 与备选方案 B 综合性能对比}
  \label{tab:q3_compare}
  \begin{tabular}{lccp{6.5cm}}
    \toprule
    \textbf{评价维度} & \textbf{方案 A (精简长效，正式采纳)} & \textbf{方案 B (多段接力)} & \textbf{工程决策权衡与对比分析} \\
    \midrule
    \textbf{中继架次数} & \textbf{3 架次} & 4 架次 & 方案 A 减少 1 次起降作业，降低起降故障风险 \\
    \textbf{保活驻留策略} & \textbf{长效静态驻留保活} & 动态分段接力 & 方案 A 避免空中二次调机，通信保障连续稳定 \\
    \textbf{系统联合总能耗} & \textbf{81.9963 kWh} & 81.8158 kWh & 能耗相差仅 0.1805 kWh（相对差异 0.22\%），方案 A 起降与交接风险更低 \\
    \textbf{系统联合完工时间} & \textbf{7746.3 s} & 7746.3 s & 均由 RT01 返航时刻决定（运输机群已于 7730.5 s 提前归巢） \\
    \textbf{实体机占用规模} & 2 架 (R01, R02) & 2 架 (R01, R02) & 均在 2 架库存限额之内，设备利用充分 \\
    \textbf{能源组件占用数} & \textbf{3 组 (MOD\_01$\sim$03)} & 4 组 & 方案 A 节约 1 组组件，保留 50\% 应急备用 \\
    \textbf{最低返航 SOC} & 30.54\% (RT01) & 64.53\% & 均满足 $\ge 20.0\%$ 安全下限，电量余量充足 \\
    \bottomrule
  \end{tabular}
\end{table}

\begin{figure}[htbp]
  \centering
  \includegraphics[width=0.65\textwidth]{fig12_q3_scheme_pareto_radar.pdf}
  \caption{中继协同调度方案 A 与方案 B 综合性能 7 维雷达决策图}
  \label{fig:q3_radar}
\end{figure}
